% !PRINTED none

\begin{center}
\gesperrt{Fifteenth Chapter.}

\textbf{The Form Problem of the Group $G_{168}$ and the Theory of Equations of the Seventh Degree.}

\rule{4cm}{0.4pt}

\S.\,125.

\gesperrt{The resolvents of the form problem.}
\end{center}

The form problem for the group $G_{168}$ (\S.\,43) yields, according to the general theory, first of all an equation of degree 168, whose coefficients are invariants. But every divisor of the group leads to a resolvent of lower degree, and indeed of degree equal to the index of the divisor.

We shall here consider only the two most interesting cases of such divisors, namely the octahedral group

\weq{1}{G_{24} = \chi^\lambda \omega^\mu \Theta^\nu,}

and the group of degree 21

\weq{2}{G_{21} = \tau^\varrho \chi^\lambda \ (\S.\,72)}

which lead us to resolvents of the seventh and eighth degrees.

As regards first the resolvents of the seventh degree, in forming them we may start from the principal axes of the group. A principal axis, namely, remains unchanged by the group $\omega^\mu \Theta^\nu$, and is carried by the whole group $G_{24}$ into three different lines. There must therefore be seven triples of principal axes, which are determined by an equation of the seventh degree.

If, following \S.\,118, (7), we set:

\weq{3}{\begin{aligned}
A_1 &= \alpha\gamma x_1 + \gamma\beta\varepsilon^3 x_2 + \alpha\beta\varepsilon^2 x_3 \\
A_2 &= \beta\gamma\varepsilon^3 x_1 + \alpha\beta\varepsilon^2 x_2 + \alpha\gamma x_3 \\
A_3 &= \alpha\beta\varepsilon^2 x_1 + \alpha\gamma x_2 + \beta\gamma\varepsilon^3 x_3,
\end{aligned}}
%CONT>>
